% Part: many-valued-logic
% Chapter: sequent-calculus
% Section: rules-and-proofs

\documentclass[../../../include/open-logic-section]{subfiles}

\begin{document}

\olfileid{mvl}{seq}{rul}

\olsection{قواعد و \usetoken{P}{derivation}}

در ادامه، بگذارید~$\Gamma, \Delta, \Pi, \Lambda$ نمایندهٔ دنباله‌های
متناهی و احتمالاً تهی از~!!{sentence}ها باشند.

\begin{defn}[دنباله]
یک \emph{دنبالهٔ $n$سویه} عبارتی به‌صورت
\[
\Gamma_1 \nSequent \dots \nSequent \Gamma_n
\]
است که در آن هر~$\Gamma_i$ دنباله‌ای متناهی و احتمالاً تهی از
!!{sentence}های زبان~$\Lang L$ است.
\end{defn}

\begin{defn}[دنبالهٔ آغازین]
یک \emph{دنبالهٔ آغازین $n$سویه}، برای هر~!!{sentence} چون~$!A$ در
زبان، دنباله‌ای~$n$سویه به‌صورت
$!A \nSequent \dots \nSequent !A$ است.

اگر زبان شامل رابط صفرموضعی~$\star$، یعنی یک ثابت گزاره‌ای، باشد،
دنباله‌ای را نیز آغازین می‌گیریم که~$\star$ دقیقاً در جایگاه متناظر با
ارزش صدق~$\tf{\star}[\Log L] \in V$ قرار دارد و همهٔ جایگاه‌های دیگرش
تهی‌اند؛ یعنی~$\dots \nSequent \star \nSequent \dots$.
\end{defn}

برای هر رابط در منطق~$n$ارزشی~$\Log{L}$، به‌ازای هر ارزش صدقی که آن
رابط می‌تواند در~$\Log{L}$ بگیرد، یک قاعدهٔ منطقی وجود دارد.
!!^{derivation}ها در حساب دنباله‌ای~$n$سویهٔ~$\Log{L}$ درخت‌هایی از
دنباله‌ها هستند که بالاترین دنباله‌هایشان دنباله‌های آغازین‌اند؛ و اگر
دنباله‌ای زیر یک یا چند دنبالهٔ دیگر قرار گیرد، باید به‌درستی با یکی از
قواعد استنتاج رابط‌های~$\Log{L}$ از آن‌ها به دست آید.

\begin{defn}[قضایا]
!!^a{sentence} چون~$!A$ یک \emph{قضیه} از منطق~$n$ارزشی~$\Log{L}$ است
اگر~!!a{derivation} از دنبالهٔ~$n$سویه‌ای وجود داشته باشد که~$!A$ را
در هر جایگاه متناظر با یک ارزش صدق ممتاز~$\Log{L}$ و دنبالهٔ تهی را در
هر جایگاه ناممتاز دارد. اگر~$!A$ قضیه باشد
$\Proves[\Log{L}] !A$، و اگر نباشد~$\Proves/[\Log{L}] !A$ می‌نویسیم.
\end{defn}

\begin{defn}[!!^{derivability}]
!!^a{sentence} چون~$!A$ را در منطق~$n$ارزشی~$\Log{L}$
\emph{!!{derivable} از} مجموعه‌ای از !!{sentence}ها چون~$\Gamma$
می‌نامیم و~$\Gamma \Proves[\Log{L}] !A$ می‌نویسیم، اگر و تنها اگر زیرمجموعهٔ متناهی
$\Gamma_0 \subseteq \Gamma$ و دنباله‌ای~$\Gamma_0'$ از
!!{sentence}های~$\Gamma_0$ وجود داشته باشد که دنبالهٔ زیر
!!a{derivation} داشته باشد:
\[ \Lambda_1 \nSequent \dots \nSequent \Lambda_n \]
که در آن اگر جایگاه~$i$ متناظر با ارزش صدقی ممتاز باشد، $\Lambda_i$
برابر~$!A$، و در غیر این صورت برابر~$\Gamma_0'$ است. اگر~$!A$ از
$\Gamma$ !!{derivable} نباشد، $\Gamma \Proves/ !A$ می‌نویسیم.
\end{defn}

برای نمونه، منطق سه‌ارزشی \L ukasiewicz یک حساب دنباله‌ای سه‌سویه دارد.
در دنبالهٔ سه‌سویهٔ~$\Gamma \nSequent \Pi \nSequent \Delta$،
$\Gamma$ متناظر با~$\False$، $\Delta$ متناظر با~$\True$، و~$\Pi$
متناظر با~$\Undef$ است. اصل‌ها به‌صورت
$!A \nSequent !A \nSequent !A$ هستند. چون فقط~$\True$ ممتاز است،
$\Gamma \Proves[\LogLuk[3]] !A$ اگر و تنها اگر دنبالهٔ
$\Gamma \nSequent \Gamma \nSequent !A$ دارای~!!a{derivation} باشد.
(اگر~$\Undef$ نیز ممتاز بود، به~!!a{derivation} از دنبالهٔ
$\Gamma \nSequent !A \nSequent !A$ نیاز داشتیم.)

\end{document}
